\documentclass[10pt,a4paper]{amsart}
\usepackage[margin=19mm]{geometry}
\usepackage{amsmath,amssymb,mathtools}
\usepackage[scr=rsfs,cal=euler]{mathalfa}
\usepackage{xfrac}
\usepackage[unicode,hidelinks,bookmarks=false]{hyperref}
\newcommand{\ie}{i.e.\ }
\newcommand{\cf}{cf.\ }
\newcommand{\resp}{respectively\ }
\newcommand{\Subsection}{\subsection}
\providecommand{\nobreakdash}{\nobreak\hbox{-}}
\setlength{\emergencystretch}{2em}
\multlinegap0pt
\usepackage{mathtools}
\usepackage{amssymb}
\usepackage{algorithm}\usepackage{algpseudocode}
\usepackage[normalem]{ulem}
\usepackage[shortlabels]{enumitem}
\newtheorem{assumption}{Assumption}
\newtheorem{theorem}{Theorem}
\newtheorem{corollary}{Corollary}
\newtheorem{proposition}{Proposition}
\usepackage{cleveref}
\crefname{assumption}{Assumption}{Assumptions}
\crefname{theorem}{Theorem}{Theorems}
\crefname{corollary}{Corollary}{Corollarys}
\crefname{proposition}{Proposition}{Propositions}
\DeclareMathOperator*{\E}{\mathbb{E}}
\title{\Large From Isotonic to Lipschitz Regression:\\ A New Interpolative Perspective on Shape-restricted Estimation}
\author{Kenta Takatsu, Tianyu Zhang, Arun Kumar Kuchibhotla}
\date{Source: arXiv:2307.05732v5 (CC BY 4.0)}
\begin{document}
\maketitle

\noindent\emph{Source and licence.} \Large From Isotonic to Lipschitz Regression:\\ A New Interpolative Perspective on Shape-restricted Estimation. Version arXiv:2307.05732v5 (CC BY 4.0), distributed by arXiv under the Creative Commons Attribution 4.0 International licence, http://creativecommons.org/licenses/by/4.0/, as recorded in the arXiv licence metadata for this exact version and shown on the abstract page https://arxiv.org/abs/2307.05732v5. Full licence notice: \texttt{licenses/arxiv-2307.05732-CC-BY-4.0.txt}.\par\smallskip

\noindent\textit{Editorial scope.} Counted unit: the article's corollary case (source0.tex lines 2541--2549). The article's complete original proof of it is delivered with it (source0.tex lines 2550--2585, one \texttt{proof} environment of 36 lines). No mathematics is written, repaired or re-derived: the statement and the proof are the article's own lines, in the article's own notation, and no retained line is substituted or re-flowed.  Retained uncounted as the source-local context the proof needs: lines 88--90; lines 330--335; lines 336--341; lines 363--368; lines 395--421; lines 425--430; lines 435--442; lines 590--596.  Nothing was omitted from any retained span, so the package carries no omission record.  No retained line contains a citation, so the wrapper carries no bibliography and no bibliography entry.  No retained line contains a figure environment, an image-inclusion command or drawing code, so no image is delivered and the package declares no asset dependency.  Editorial formatting only, in addition: an AMS \texttt{amsart} preamble that restates the article's own declarations and macros used by the retained lines, namely the environments \texttt{assumption}, \texttt{theorem}, \texttt{corollary}, \texttt{proposition} on separate counters, together with the standard packages those lines need.  The retained environments print in this document as Assumption 1, Theorem 1, Corollary 1, Proposition 1 in order of appearance, whereas the article prints the same items with the numbering of its own full document.  The complete native source of the article as distributed in the arXiv e-print for this exact version is bundled unchanged as \texttt{sources/144156/source0.tex}.
\begin{align}\label{eq:regression-function}
    Y_i = f_0(X_i) + \xi_i \, \text{ for }\, i = 1, \dots, n.
\end{align}

\begin{assumption}\label{as:q-moment}
    There exists constants $q \geq 2$ and $C_q \in (0,\infty)$ such that 
    \begin{equation*}
        \left(\E[|\xi_i|^q\mid X_i]\right)^{1/q}\le C_q \text{ almost surely}.
    \end{equation*}
\end{assumption}

\begin{assumption}\label{as:beta-weibull}
    There exist constants $\beta>0$ and $C_\beta \in(0, \infty)$ such that for all $r \geq 1$
    \begin{equation*}
\left(\mathbb{E}\left[|\xi_i|^r\mid X_i\right]\right)^{1 / r} \leq C_\beta r^{1 / \beta}\text{ almost surely}.
\end{equation*}
\end{assumption}

    \begin{align}\label{eq:rate_under_LSE}
    	\mathrm{R}_n := \begin{cases}
    		n^{-1+1/q} &\text{ under }\Cref{as:q-moment},\\
    		n^{-1}(\log n)^{1/\beta} &\text{ under }\Cref{as:beta-weibull}.
    	\end{cases}
    \end{align}

\begin{theorem}\label{th: general_isotonic}
Consider the IID setting \eqref{eq:regression-function}.
Suppose one of the Assumptions \ref{as:q-moment} or \ref{as:beta-weibull} hold. Then, there exists a constant $N_\epsilon$ only depending on $\epsilon\in(0,1)$ such that for any $n \ge N_\epsilon$, with probability greater than $1-\epsilon$,
\begin{equation*}
\begin{aligned}
\|\widehat{f}_n-f_0\|^2  \leq \inf _{L \in \mathcal{L}}\inf _{1 \leq m \leq n}\left\{ \, \inf _{f \in \mathcal{F}_m(1, L)}\left\|f-f_0\right\|^2+ 
\mathrm{R}_{L,m}^{\mathrm{Est}}\right\} +\mathrm{R}^{\rm CV}_{\rm iso},
\end{aligned}
\end{equation*}
where
\begin{equation*}
    \begin{aligned}
\mathrm{R}_{L,m}^{\mathrm{Est}}& =C_\epsilon m \sigma^2 B_n^2 \log ^2\left(n \right)n^{-1},\\
        \mathrm{R}^{\rm CV}_{\rm iso}
 & =  C_\epsilon \log (1+ L_{+} n^{1 / 2})\left( B_n\mathrm{R}_n + B_n^2n^{-1} \right), 
 \end{aligned}
 \end{equation*}
 and
\begin{equation*}
\mathrm{R}_n =
\begin{cases}
   n^{-1+1/q}, & \text{under } \Cref{as:q-moment}, \\[6pt]
   n^{-1} (\log n)^{1/\beta}, & \text{under } \Cref{as:beta-weibull}.
\end{cases}
\end{equation*}
We used $B_n=\left\|f_0\right\|_{\infty}+\sigma^2+L_{+}$, which may depend on $n$ though $L_+$.
\end{theorem}

 \begin{corollary}[Low complexity adaptation with \texttt{MEAN}]\label{cor:low-complexity} Consider the IID setting \eqref{eq:regression-function} and assume $f_0 \in \mathcal{F}_m(1, L_0)$ for $L_0 \ge 0$. We assume $L_+\rightarrow \infty$ as $n \rightarrow\infty$ and one of the noise conditions \ref{as:q-moment} or \ref{as:beta-weibull} holds. Then, for any $\epsilon \in(0,1)$ and large enough $n$, we have
\begin{equation*}
\|\widehat{f}_n-f_0\|^2 \leq C_\epsilon\left(\frac{m L_{+}^2 \log ^2 n}{n}+\log \left(1+L_{+} n\right)\left(L_{+} \mathrm{R}_n+L_{+}^2 n^{-1}\right)\right),
\end{equation*}
with probability greater than $1-\epsilon$. Here, $C_\epsilon > 0$ is a constant depending on $\epsilon, L_0, \|f_0\|_\infty, \sigma^2$ as well as constants in \ref{as:q-moment} or \ref{as:beta-weibull}, and $\mathrm{R}_n$ is defined in \Cref{eq:rate_under_LSE}.
\end{corollary}

\begin{corollary}[Worst case guarantee with \texttt{MEAN}]\label{cor:worst-case} 
Assume the same setting as in \Cref{cor:low-complexity}, but with $f_0\in\mathcal{F}(1,L_0)$ for $L_0 \ge 0$.
Then, for any $\epsilon\in(0,1)$ and large enough $n$, we have
\begin{equation*}
\|\widehat{f}_n-f_0\|^2 \leq C_\epsilon\left(\frac{L_{+}^{4 / 3}(\log n)^{4 / 3}}{n^{2 / 3}}+\log \left(1+L_{+} n\right)\left(L_{+} \mathrm{R}_n+L_{+}^2 n^{-1}\right)\right),
\end{equation*}
with probability greater than $1-\epsilon$. Here, $C_\epsilon >0$ is a constant depending on $\epsilon, L_0, \|f_0\|_\infty, \sigma^2$ as well as constants in \ref{as:q-moment} or \ref{as:beta-weibull}, and $\mathrm{R}_n$ is defined in \Cref{eq:rate_under_LSE}.
\end{corollary}

\begin{proposition}\label{prop: BV approximation}
For any function $f_0 \in \operatorname{BV}(1, M;[a,b])$, there exists a function $f_{L}\in\mathrm{BL}_1(1, L ; \Omega)\subset\mathcal{F}^\dagger(1,L;[a,b])$ such that 
\begin{equation*}
\int_a^b\left(f_L(x)-f_0(x)\right)^2 d x \leq CM^3 L^{-1},
\end{equation*}
where $C>0$ is a universal constant.
\end{proposition}

\begin{corollary}[BV truth for \texttt{MEAN}]\label{cor: BV case} 
Assume the same setting as in \Cref{cor:low-complexity}, except that $f_0 \in \operatorname{BV}(1, M; [0,1])$. Additionally, suppose that $P_X$ admits a bounded density with respect to the Lebesgue measure on $[0,1]$. Let $\mathcal{L} = [0,n^{2/9} \log^{-4/9} n]$.

Then, for any $\epsilon\in(0,1)$ and large enough $n$, we have
\begin{equation*}
\left\|\widehat{f}_n-f_0\right\|^2 \leq C_\epsilon\left(\frac{\log ^{4 / 9} n}{n^{2 / 9}}+\log ^{5 / 9} n \cdot n^{2 / 9} \mathrm{R}_n\right)
\end{equation*}
with probability greater than $1-\epsilon$. Here, $C_\epsilon \in (0, \infty)$ is a constant depending on $\epsilon, \|f_0\|_\infty, \sigma^2$ as well as constants in \ref{as:q-moment} or \ref{as:beta-weibull}, and $\mathrm{R}_n$ is defined in \Cref{eq:rate_under_LSE}.
\end{corollary}

\begin{proof}
The proof of \Cref{cor: BV case} is similar to that of \Cref{cor:worst-case}, both need \Cref{th: general_isotonic}:
\begin{equation*}
\left\|\widehat{f}_n-f_0\right\|^2 \leq \inf _{L \in \mathcal{L}} \inf _{1 \leq m \leq n}\left\{\inf _{f \in \mathcal{F}_m(1, L)}\left\|f-f_0\right\|^2+\mathrm{R}_{L,m}^{\mathrm{Est}}\right\}+\mathrm{R}_{\mathrm{iso}}^{\mathrm{CV}}.
\end{equation*}
Consider an $l_n$-Lipschitz function $f_{l_n}$, we bound $\left\|f-f_0\right\|^2$ by
\begin{equation*}
\left\|f-f_0\right\|^2 \lesssim\left\|f-f_{l_n}\right\|^2+\left\|f_{l_n}-f_0\right\|^2.
\end{equation*}
Then we have 
\begin{equation}\label{eq: bv before plug in}
\left\|\widehat{f}_n-f_0\right\|^2 \lesssim \inf _{L \in \mathcal{L}} \inf _{1 \leq m \leq n}\left\{\inf _{f \in \mathcal{F}_m(1, L)}\left\|f-f_{l_n}\right\|^2+\mathrm{R}_{L,m}^{\mathrm{Est}}\right\}+\mathrm{R}_{\mathrm{iso}}^{\mathrm{CV}}+\left\|f_{l_n}-f_0\right\|^2.
\end{equation}
In the proof of \Cref{cor:worst-case} we have established that 
\begin{equation*}
\left\|f-f_{l_n}\right\|^2 \leq\left(\left\|f_{l_n}\right\|_{\infty}^2+l_n\right)^2 m^{-2}.
\end{equation*}
Applying \Cref{prop: BV approximation}, we also know that 
\begin{equation*}
\left\|f_{l_n}-f_0\right\|^2 \leq C\left(f_0, P_X\right) l_n^{-1}
\end{equation*}
for $P_X$ with a bounded density function.

Plug these quantities into \Cref{eq: bv before plug in} (omitting $\mathrm{R}_{\text {iso }}^{\mathrm{CV}}$ for now), we need to find $m,l_n$ that minimizes
\begin{equation*}
\left(\left\|f_{l_n}\right\|_{\infty}^2+l_n\right)^2 m^{-2}+C_\epsilon m \sigma^2 B_n^2 \log ^2(n) n^{-1}+C\left(f_0, P_X\right) l_n^{-1} .
\end{equation*}

One should use $m = n^{1/3}$, $L_+ = l_n = n^{2/9}\log^{-4/9} n$ to balance the three terms above. These choices give
\begin{equation*}
\begin{aligned}
\left\|\widehat{f}_n-f_0\right\|^2 & \lesssim C_\epsilon \log ^{4 / 9} n \cdot n^{-2 / 9}+\mathrm{R}_{\mathrm{iso}}^{\mathrm{CV}} \\
& \lesssim C_\epsilon\left(\log ^{4 / 9} n \cdot n^{-2 / 9}+\log ^{5 / 9} n \cdot n^{2 / 9} \mathrm{R}_n\right).
\end{aligned}
\end{equation*}
\end{proof}

\end{document}
